% !TEX root = micro_Lie_theory.tex

%%%%%%%%%%%%%%%%%%%% S3 cont %%%%%%%%%%%%%%%%%%%%%%%%%%%%%%%%%
\begin{fexample}{The unit quaternions group $S^3$ (cont.)}
\label{ex:S3}
In the  group $S^3$ (recall \exRef{ex:S3_intro} and see \eg\ \cite{SOLA-17-Quaternion}), the time derivative of the unit norm condition $\bfq^*\bfq=1$ yields 
%
$$\bfq^*\dot\bfq=-(\bfq^*\dot\bfq)^*
.
$$
%
This reveals that $\bfq^*\dot\bfq$ is a pure quaternion (its real part is zero). 
%The set of pure quaternions is noted $\bbH_p$.
Pure quaternions $\bfu v\in\bbH_p$ have the form 
%
$$\bfu v=(iu_x+ju_y+ku_z)v =iv_x+jv_y+kv_z 
,$$
%
where $\bfu\te iu_x+ju_y+ku_z$ is pure and unitary, $v$ is the norm, and $i,j,k$ are the generators of the Lie algebra $\frak{s}^3=\bbH_p$.
%
Re-writing the condition above we have,
%
\begin{align*}
\dot\bfq=\bfq\, \bfu v &&\in \mtanat{S^3}{\bfq}
,
\end{align*}
% 
which integrates to $\bfq = \bfq_0\exp(\bfu v t)$. 
Letting $\bfq_0=1$ and defining $\bphi \te \bfu\phi \te\bfu v t$ we get the exponential map,
%
\begin{align*}
\bfq = \exp(\bfu\phi) \te \sum \frac{\phi^k}{k!}\bfu^k &&\in S^3
~.
\end{align*}
%
The powers of $\bfu$ follow the pattern $1,\bfu,-1,-\bfu,1,\cdots$.
Thus we group the terms in $1$ and $\bfu$ and identify the series of $\cos\phi$ and $\sin\phi$.
We get the closed form,
%
\begin{align*}
\bfq = \exp(\bfu\phi) = \cos(\phi) + \bfu\sin(\phi) 
~,
\end{align*}
%
which is a beautiful extension of the Euler formula, $\exp(i\phi)=\cos\phi+i\sin\phi$.
%
The elements of the Lie algebra $\bfphi=\bfu \phi\in\frak{s}^3$ can be identified with the rotation vector $\bth\in\bbR^3$ trough the mappings \emph{hat} and \emph{vee},
%
\begin{align*}
\textrm{Hat}&:& \bbR^3&\to\frak{s}^3;& \bth&\mapsto\bth^\wedge=\bth/2
\\
\textrm{Vee}&:& \frak{s}^3&\to\bbR^3;& \bphi&\mapsto\bphi^\vee=2\bphi
~,
\end{align*}
%
where %$\bth \te \bfu\theta\te\bw t$ is the angle-axis vector, and 
the factor 2 accounts for the double effect of the quaternion in the rotation action, $\bfx'=\bfq\,\bfx\,\bfq^*$. 
%
With this choice of Hat and Vee, the quaternion exponential
%
\begin{align*}
\bfq = \Exp(\bfu\theta)=\cos(\theta/2)+\bfu\sin(\theta/2)
\end{align*}
%
is equivalent to the rotation matrix $\bfR=\Exp(\bfu\theta)$.
\end{fexample}
%%%%%%%%%%%%%%%%%%%%%%%%%%%%%%%%%%%%%%%%%%%%%%%%%%%%%%%%
